\chapter[Fundamentação Teórica]{Fundamentação Teórica}

\section{Blockchain}

Uma blockchain é uma rede \textit{p2p} que compartilha um livro-razão
(ledger) imutável, atualizável apenas por meio de consenso ou acordo
entre os nós participantes da rede \cite{MasteringBlockchain}. Em suma,
uma blockchain é um banco de dados distribuído que armazena um conjunto de
transações, imutáveis (não podem ser alteradas ou apagadas) e que apenas
são adicionadas quando há quórum entre os participantes da rede.

A tecnologia foi proposta por Satoshi Nakamoto em seu artigo Bitcoin:
A Peer-to-Peer Electronic Cash System, como uma solução criar um
sistema de pagamentos on-line que não dependa de intermediários
financeiros para validar transações \cite{Bitcoin}.

Essas características fazem das blockchains uma tecnologia ideal para
aplicações que exigem transparência, segurança e descentralização, como
sistemas financeiros, cadeias de suprimentos, votação eletrônica, entre
outros.

\subsection{Mecanismos de consenso}

O conjunto de regras e procedimentos que permitem que os nós de uma
rede cheguem a um acordo sobre o estado atual do livro-razão é
conhecido como mecanismo de consenso. \cite{MasteringBlockchain}.

Existem diversos mecanimos de consenso, cada qual com suas vantagens
e desvantagens.

Dentre os mais conhecidos, destacam-se:

\subsubsection{Prova de Trabalho (Proof of Work - PoW)}

Como o nome sugere, o mecanismo de Prova de Trabalho (PoW) exige que os
nós participantes da rede realizem um trabalho computacional para validar
transações e adicionar novos blocos à rede.

No Bitcoin, o trabalho consiste em encontrar um valor chamado nonce
que, ao ser combinado com os dados de um conjuto de transações
(bloco) e processado pela função
hash SHA-256, produza um hash que satisfaça a seguinte condição:

$$
hash < target
$$

Onde $hash$ é o resultado do algortimo SHA-256 aplicado ao bloco e
$target$ é um valor númerico definido pela rede que determina a
dificuldade do problema.
\cite{EthereumWhitepaper}

Ao resolver o problema, o bloco é enviado para os demais nós da rede.
Cada nós valida o bloco e o adiciona à sua cópia do livro-razão, caso
o bloco seja considerado válido. O nó que resolver o problema
primeiro é recompensado com uma quantia da criptomoeda da rede, além
das taxas de transação associadas ao bloco.

O mecanismo de PoW é um mecânismo robusto e seguro, pois é inviável
econômicamente para um nó malicioso atacar a rede, uma vez que,
exigiria um alto poder computacional para superar capacidade de
processamento da rede.
Além disso, seria mais rentável para o nó malicioso participar
honestamente da rede
do que tentar atacá-la.

Entretanto, o PoW é criticado por seu alto consumo energético, uma vez
que a resolução dos problemas computacionais exige um grande poder
de processamento, o que leva a um consumo elevado de energia elétrica.
\cite{BitcoinEnergyConsumption}

\subsubsection{Prova de Participação (Proof of Stake - PoS)}

Em uma rede que utiliza o mecanismo de Prova de Particiapação (PoS),
os nós validadores
precisam depositar uma quantia da criptpmoeda da rede como garantia
(\textit{stake})
para participar do processo de validação. Caso o nó atue de forma
maliciosa, o valor depositado
é perdido. \cite{ethereum-docs}

Apesar de ser um mecanismo mais eficiente em termos energéticos do que o PoW,
o PoS apresenta desafios relacionados à segurança e à centralização,
uma vez que nós com grandes quantias de criptomoeda podem ter mais
influência na validação de blocos.

\subsubsection{Prova de Autoridade (Proof of Authority - PoA)}

A prova de autoridade é um mecanismo de consenso mais centralizado,
neste, um conjunto de nós confiáveis é
pré-selecionado para validar transações e criar novos blocos. Esses
nós são conhecidos como validadores
ou autoridades, e sua identidade é conhecida e confiável pela rede.
\cite{ethereum-docs}

\section{Criptografia}

A criptografia é frequentemente utilizada em muitos aspectos da
computação moderna. Ela é usada para enviar
mensagens de forma segura através de uma rede, bem como para proteger
dados de banco de dados,
arquivos e até mesmo discos inteiros contra a leitura de seu conteúdo
por entidades não autorizadas.
\cite{OperatingSystemsConcepts}

Um algoritmo de criptografia deve possuir duas funções principais:

\begin{itemize}
  \item Função criptográfica: $E_k(m) = c$, que dado uma mensagem $m$
    e uma chave $k$, produza um texto cifrado $c$.
  \item Função de descriptografia: $D_k(c) = m$, que dado um texto
    cifrado $c$ e uma chave $k$, produza a mensagem original $m$.
\end{itemize}

A chave $k$ não necessariamente precisa ser a mesma para as funções
de criptografia $E$ e de descriptografia $D$. Essa característica é
base para a classificação dos algoritmos de criptografia em dois
tipos: criptografia de chave simétrica e criptografia de chave assimétrica.

\subsection {Criptografia de chave assimétrica}

\subsection{Critografia de chave simétrica}

\section{Funções hash}

\section{Lei Geral de Proteção de Dados Pessoais (LGPD)}

A Lei Geral de Proteção de dados Pessoais (LGPD) dispõe sobre o
tratamento de dados pessoais, com o objetivo de proteger os direitos
fundamentais de liberdade e de privacidade e o livre desenvolvimento
da personalidade da pessoa natural. \cite{LGPD}

Uma das principais implicações da LGPD para o segmento da saúde se
refere à forma como a coleta,
o armazenamento e todo o tratamento de dados serão empregados. A lei
exige o consentimento do paciente para
uso de seus dados pessoais para uma finalidade determinada.
\cite{LGPD-Saude}.

Ainda assim existe a possibilidade de tratamento de dados pessoais
sem o consentimento do paciente, desde que seja para finalidades
específicas, como cumprimento de obrigação legal ou regulatória,
execução de políticas públicas, realização de estudos por órgão de
pesquisa, entre outros. \cite{LGPD-Saude}.
